%% AUTO-GENERATED by `python -m experiments.report` -- do not hand-edit.
%% Schema: model x device; n=1000 timed runs after 50 warm-up, batch 1, fp32, against the 4 ms GOOSE transfer-time envelope
%% Every cell comes from a RunResult; a red numeric marker flags a value the
%% loaded runs cannot support, n.a. marks a structurally inapplicable cell.
%% Layout normalized to the group table template (tabular* full width, scriptsize, \tablenota).
\scriptsize
\setlength{\tabcolsep}{2pt}%
\renewcommand{\arraystretch}{1.15}%
\begin{tabular*}{\linewidth}{@{\extracolsep{\fill}}l l c c c c c@{}}
\toprule
Model & Device & Mean (ms) & p50 (ms) & p95 (ms) & p99 (ms) & Share of 4\,ms env.\ (\%) \\
\midrule
CNN-LSTM & CPU & 1.108 & 0.678 & 2.878 & 8.486 & 212.1 \\
CNN-LSTM & GPU & 1.013 & 0.729 & 2.412 & 3.280 & 82.0 \\
\midrule
Spiking & CPU & 5.762 & 4.940 & 10.519 & 16.070 & 401.8 \\
Spiking & GPU & 8.303 & 7.821 & 11.920 & 16.934 & 423.3 \\
\midrule
Mamba & CPU & 10.552 & 9.430 & 17.034 & 24.582 & 614.5 \\
Mamba & GPU & 11.307 & 10.649 & 16.400 & 21.907 & 547.7 \\
\bottomrule
\end{tabular*}
\tablenota{Exact cells of Figure~\ref*{fig:latencyenv}. fp32 only: torch's dynamic INT8 kernels reject the recurrent and state-space layers. One fixed reference machine, 50 warm-up runs; values bound the compute share of the 4\,ms envelope, not compliance.}